\documentclass[../Main.tex]{subfiles}

\begin{document}

\section{Reflexion} \label{sec:Reflexion}

\subsection{Reflexion des Arbeitsprozesses}

Dieses projekt hat etwa 1 jahr unserer zeit in anspruch genommen und ist durch verschiedene pahsen durchgegangen hier möchten wir diese der reihe nach aufgreifeun und unsere gedanken dazu abgeben.

ganz am anfang unserer arbeit, bei der wir uns noch völlig unsier waren und völlig andere vorstellungen hatten, was am ende für ein produkt dabei entstehen würde haben wir mal in alle berreiche etwas reingestochert, wir haben mit der gamen spielesystemen angefangen, von denen wir noch ohne recherche ausggegenagen sind, dass wir sie brauchen, wir haben angefangen uns zu überlegen wei wir unsere forschungsfrage genau beantworten wollen und wir haben damit angefangen zu recherchieren wie man recherchiert. vor uns waren noch fast ein jahr zeit und die erste deadline die in 4-5-6 monaten fällig war(bei zeitangaben musst du prüfen). wir sind noch wenig zielorientiert und machen gerade einfach das was uns gefällt und es wirkt vielleicht etwsa planlos ist es aber rückbetrachtend nicht. wir haben einfach mal ausgekundschaftet was es gibt. die fallhöhe war hier allerdings auch noch nicht so how weil egal was wir gemacht hätten hätte nichts wirklich verworfen werden müssen und die prozesse hingen jetzt noch nicht wirklich voneinandner ab. es war eher ein glücklicher zufall als strukturiertes denken unserer seits aber die natur ein projekt anzufangen ist halt so.. und es gibt ja diesen einen spruch unter leuten die projekte machen, dass sie nie angefangen hätten, hätetn sie gewusst, dass es so viel wäre. könnte ich es noch einmal mahen, würde ich es genau so wieder angehen, mit dem einzigen nterschied, dass ich wüsste dass ich die narrenfrist habe.

vertraglich hatten wir mittlerweile alles mit unserer Betreuungsperson festgehalten, was wir mussten. und bis zur ersten deadline spielte die beantwortung der der erkenntnisfrage eher eine untergeordnete rolle. das lag besonders daran, dass für die playtest und damit auch für die erkenntnisfrage relevanten systeme auf einer rechnischen basis beruhten, die wir bis zu diesem zeitpunkt noch nicht hatten. also arbeiteten wir weiter an diesen vorraussetzungen. heisst einen character mit allen nötigen animationen, code der diesen spieler steuert, gegner code, code das überhaupt ein game zum starten bringt. diese zeit würde ich rückblickent auch von uns ziemlich gut entshcieden ansehen. obwohl wir etwas den recherche teil vernachlässigt haben und am prototyp gearbeitet haben, war das die richtige entschidung. wir haben das zum einen gemacht, weil wir an diesem teil am meisten freude fanden aber auch weil wir unsere orstellung von unserem eigenen spiel bilden mssten. dabei war gegenseitige kommuniktion besonders wichtig. nachdem wir shcon die ersten systeme erstellt hatten, hatten wir gelernt, worauf wir in anderen spielen achten mussten und mit diesem wissen konnen wir erst recht eine gute recherche machen. könnten wir das noch einmal so machen, würden wir es abhängig von unseren vorstellugnen und wie konkret sie schon sind also je konkreter ist besser, und unserem kenntnisstand entweder zuerst die rcherche machen, wenn wir beides haben und wenn wir eines von beiden noch nicht haben, dann wieder zuerst eigene erfarungen sammeln

gleichzeitig rückte die erste deadline die Zwischenpräsentation näher. heisst wir mussten zuschauen, dass wir unssee unterilegenden systeme fertig bekommen. das hiess jetzt schon, das wir features die wir wollten aber nicht besonders relevant waren für unsere erkenntnisfrage schon streichen mussten. oder stark eingrenzen und wir mussten das priorisiren, was für ein spiel essenziell ist. dazu kam noch das wir ein palytest gemacht haben mussten, welches auf einer begründeten basis gemacht wurde. in dieser zeit hat auch besonders zeitmanagement eine grosse rolle gespielt. mit dem vertrag hatten wir eine checkliste, die noch viele fehlende teile hatte. bis zur abgabe haben wir allerdings alles hinbekommen, auch wenn es sehr sehr sehr knapp wurde. nicht das vertigstellen sondern besonders die qualitätssicherung hat uns hier einen strich durch die rechnung gemahct und da wir für uns so wenig kompromisse eingehen wollten und möglichst viel qualität wollten, ist das so gelaufen. im nachhinein wäre es auch interessant gewesen, zu wissen wie nicht 100 prozent sondern vielleicht nur 80 oder 85 prozent leistund und qualität bewertet wurden. da auch zwischenbwertung nicht gezählt hat, ist das eine verpasste chance. anderer seit haben wir die erfahrung mit crunchen machen können die uns am ende die fähigkeit gegeben hat, 95 bis 100 prozent rauszuholen mit genügent zeit und kaum crunchen. das wird später noch bei der finalen abgabe thematisiert.

zur zwischen abgabe haben wir dann alles was wir vertraglich festgehalten haben erreicht. und dazu haben wir natürlich auch ein feedback bekommen. vieles war gut einiges auch weniger. die kritik die wir erhalten haben betraf weniger den prototyp sondern mehr alles rund herum. besonders die begleitenden schriften hatten einige schwachstellen. wir haben versucht diese kritik umzusetzen, wir hatten zunächst jedoch schwierikkeiten dabei, da wir die angeprangerten punkte nich t verinnerlicht hatten, sonst hätten wir sie zu begin gleich richtig gemacht. deshalbt stellten wir unserer begleitperson einige fragen, unzwar so lange bis wir uns sicher waren, es verstanden zu haben.

nach der zwischenpräsentation hatten wir wieder genug zeit und die deadline war genug weit weg, dass wir unser escope ambizionen vergrösserten. gleichzeitig nahm die anfangsgescwhindigkeit, mit der wir neues produzierten ab das lag zu einen daran, dass das projekt nun so gross war, dass kleinere neuheiten wenig unterschied machen. und der anfängliche eifer ist zwar weit davon erloschen zu sein aber das projekt zu maintaien und auch sachen zu machen, die uns weniger spass machen und die sich mitlerweile angehäuft haben setzt uns langsam zu

\begin{itemize}[leftmargin=1.5cm]
    \item \textcolor{red}{\textbf{1. Ausgangslage und Plan:} Erste Commits im Juni 2025; der Vertrag sah vor: Zwischenprodukt 10.6., Zwischenpräsentation 17.6., Analyse der Referenzspiele Anfang Juli, Spielsysteme Mitte August, Doku und Fragebögen Ende September, Playtests Mitte Oktober, Verbesserungen Ende Oktober.}
    \item \textcolor{red}{\textbf{2. Verlauf in Phasen} (aus Git und Journal): Hauptentwicklung des Prototyps Februar bis Juni 2026; Juni: README und Architekturdiagramm fürs Zwischenprodukt; 16.--20.7.: Kapitel 1--4 und die \game{NieR}-Analyse als Entwurf; Juli bis September kaum Arbeit am Spiel (Sommerpause?); Oktober: Level überarbeitet, Playtests vorbereitet und durchgeführt; November: Begleitbericht fertiggestellt.}
    \item \textcolor{red}{\textbf{3. Scope-Entscheide:} 20.5.: Vertrag bereinigt, Projektumfang reduziert, Erkenntnisfrage festgelegt -- was wurde gestrichen? Später: Bonus-Inhalte weggelassen, technische Dokumentation ausgelagert, README gestrichen, Recherche in Kapitel~\ref{sec:Leveldesign} auf Rat der Betreuungsperson auf das für die Erkenntnisfrage Nötige gekürzt.}
    \item \textcolor{red}{\textbf{4. Plan und tatsächlicher Verlauf:} Playtests 26.10.--2.11.\ statt Mitte Oktober, Verbesserungen erst im Bericht statt Ende Oktober. Warum?}
    \item \textcolor{red}{\textbf{5. Wichtige Wendepunkte} (Journal 24./25.10., 13~h): Obwohl der Prototyp technisch fertig war, machten die Kämpfe bei eigenen Tests wenig Spass. Ihr habt die Encounters gekürzt, die Gegner aggressiver gemacht und Rechercheergebnisse nachgeholt, die noch nicht umgesetzt waren. Eure eigene Erkenntnis dort: \enquote{Die technische Fertigstellung allein garantiert noch kein gutes Spielerlebnis} -- und weil ihr das Level so oft gespielt habt, braucht es externe Playtests. Das ist ein starkes Reflexionsbeispiel.}
    \item \textcolor{red}{\textbf{6. Rückmeldungen der Zwischenbewertung und was ihr daraus gemacht habt:} Zeitaufwand konsequenter dokumentieren, Reflexionen vertiefen, Entscheidungen mit Alternativen begründen, KI vollständig dokumentieren, Planung mit Arbeitspaketen, Bericht und Code verknüpfen.}
    \item \textcolor{red}{\textbf{7. Was gut funktioniert hat:} z.\,B. iteratives Vorgehen, Git und GitHub, Playtest-Vorgehen -- jeweils mit Grund.}
    \item \textcolor{red}{\textbf{8. Methodik im Rückblick:} kurz, Einzelheiten in den Limitationen (Abschnitt~\ref{sec:Limitationen}).}
    \item \textcolor{red}{\textbf{9. Was ihr anders machen würdet}, z.\,B. früher extern testen, Rechercheergebnisse von Anfang an mit dem Level abgleichen, Bericht laufend statt am Schluss schreiben.}
    \item \textcolor{red}{\textbf{10. Verweis aufs Arbeitsjournal}, wo der Verlauf im Detail dokumentiert ist.}
\end{itemize}

\subsection{Gruppenarbeit und Aufgabenteilung}

\begin{itemize}[leftmargin=1.5cm]
    \item \textcolor{red}{\textbf{1. Wie ihr die Arbeit aufgeteilt habt.} Laut Zwischenbewertung: Rodin stärker bei Animation, UI, Bericht und README; Zhile stärker bei Architektur, Bewegung, State Machines und Kampf-Grundlagen. Das deckt sich mit Git: Zhile vor allem \code{Entities/Player}, \code{Common} (Utility, Tools), \code{Core}, Kamera; Rodin vor allem \code{Interfaces} und Menüs sowie die TrenchBroom-Karten des Levels.}
    \item \textcolor{red}{\textbf{2. Tabelle Aufgabenteilung und Eigenleistung} für Prototyp und Bericht (Rodin / Zhile / gemeinsam). Zum Bericht: Kapitel 1--8 und die \game{NieR}-Analyse Rodin, die \game{Ninja-Gaiden}-Analyse Zhile, README und Architekturdiagramm (Juni) vor allem Zhile. Bitte prüfen und ergänzen (Playtests, Leveldesign, Audio, \ldots).}
    \item \textcolor{red}{\textbf{3. Zusammenarbeit und Kommunikation:} Wie habt ihr euch abgesprochen (Treffen, Discord, GitHub)? Wie habt ihr Entscheidungen getroffen?}
    \item \textcolor{red}{\textbf{4. Veränderung im Verlauf:} z.\,B. ab Oktober Rodin vor allem Level, Playtests und Bericht, Zhile weiter am Spiel (laut Git im Oktober sehr aktiv). Wie kam es dazu?}
    \item \textcolor{red}{\textbf{5. Fazit:} Was lief gut, was war schwierig, was würdet ihr in einer Gruppenarbeit anders machen? Hinweis: Commit-Zahlen sagen wenig über den Aufwand aus, deshalb nicht als Massstab verwenden.}
\end{itemize}

\subsection{Verwendung von KI} \label{sec:ReflexionKI}

\begin{itemize}[leftmargin=1.5cm]
    \item \textcolor{red}{\textbf{1. Welche Tools wofür:} Code (Treppen-Werkzeug und Animations-Hilfsskript im Frühling, Audiosystem im Oktober), Begleitbericht (Korrektur, Struktur, Rückmeldungen). Details in der KI-Deklaration (Abschnitt~\ref{sec:KIDeklaration}).}
    \item \textcolor{red}{\textbf{2. Was geholfen hat:} z.\,B. Widersprüche zwischen Kapiteln finden, unbelegte Aussagen markieren, \LaTeX-Formatierung, Fehlersuche im Code.}
    \item \textcolor{red}{\textbf{3. Grenzen und Probleme:} Die KI ging mehrmals von falschen Annahmen aus, die ihr korrigieren musstet, z.\,B. dass Themenfindung mit KI gemacht wurde, dass ein Tutorial geplant war, welche Räume Brückenräume sind, oder dass in den Protokollen ein anderer Archer gemeint war.}
    \item \textcolor{red}{\textbf{4. Wie ihr die Ergebnisse überprüft habt:} Abgleich mit Quellen, Playtest-Protokollen und dem Spiel selbst; Vorschläge nicht ungeprüft übernommen, einige bewusst verworfen.}
    \item \textcolor{red}{\textbf{5. Eigenanteil sichern:} Interpretation und Urteile selbst festgelegt; von der KI formulierte Stellen markiert und in eigenen Worten neu geschrieben; Regel \enquote{nicht wortwörtlich übernehmen} aus dem Merkblatt.}
    \item \textcolor{red}{\textbf{6. Fazit:} Würdet ihr KI wieder so einsetzen? Was würdet ihr anders machen, z.\,B. KI-Nutzung von Anfang an laufend im Journal dokumentieren?}
\end{itemize}

\subsection{Persönliche Erkenntnisse}

\subsubsection{Rodin Arslan}

\begin{itemize}[leftmargin=1.5cm]
    \item \textcolor{red}{\textbf{1. Fachlich:} z.\,B. Animation (u.\,a. Hängeanimation), UI und Menüs, Leveldesign mit TrenchBroom, Gegner- und Encounter-Design, Playtesting, \LaTeX\ und BibTeX.}
    \item \textcolor{red}{\textbf{2. Methodisch:} Spielanalyse mit dem MDA-Framework, Recherche und Zitieren, Playtests planen und auswerten, Projektplanung.}
    \item \textcolor{red}{\textbf{3. Persönlich:} z.\,B. dass ein technisch fertiges Spiel noch nicht gut spielbar ist (Journal 24./25.10.), Zeitmanagement in der Schlussphase, Umgang mit Rückmeldungen.}
    \item \textcolor{red}{\textbf{4. Was du mitnimmst} für die Zukunft.}
\end{itemize}

\subsubsection{Zhile Zhu}

\textcolor{red}{ZHILE DAS IST FÜR DICH}

\end{document}
